\chapter{2009:Venkatraman Ramakrishnan and Thomas A. Steitz and Ada E. Yonath}

\label{chap:chemistry2009}

The Nobel Prize in Chemistry for the year 2009 was awarded jointly to Venkatraman Ramakrishnan and Thomas A. Steitz and Ada E. Yonath.

\textbf{Motivation:}

for studies of the structure and function of the ribosome

\section{Venkatraman Ramakrishnan}

\subsection*{Early Life and Education}

Venkatraman Ramakrishnan was born in 1952 in Chidambaram, India.

\subsection*{Career and Major Research}

Ramakrishnan earned a bachelor's degree from the Maharaja Sayajirao University of Baroda and a doctorate in physics from the University of California, San Diego, in 1976. After work at the Brookhaven National Laboratory he became a group leader at the MRC Laboratory of Molecular Biology in Cambridge, where he determined the atomic structure of the small (30S) subunit of the ribosome.

\subsection*{Nobel Prize-winning Work}

The work for which Venkatraman Ramakrishnan was awarded the Nobel Prize in Chemistry in 2009 was described by the Nobel Committee as: "for studies of the structure and function of the ribosome".

Ramakrishnan's structure of the 30S subunit showed how the ribosome reads the genetic message and how antibiotics interfere with that process.

\subsection*{Significance and Impact}

The structure explained how the small subunit ensures the accuracy of protein synthesis and revealed the binding sites of several antibiotics, guiding the design of improved drugs.

\subsection*{Later Career and Honors}

Ramakrishnan later served as president of the Royal Society from 2015 to 2020 and continues his research in Cambridge.

\subsection*{Legacy}

His work is central to the modern understanding of how the ribosome decodes messenger RNA and makes proteins.

\subsection*{Modern Developments}

Ramakrishnan's X-ray crystallographic structures of the ribosome, together with Steitz and Yonath, revealed the atomic architecture of the machine that translates mRNA into protein. His structures explain how antibiotics target the ribosome, guiding the design of new antibacterial drugs, and his work underpins modern understanding of protein synthesis.

\subsection*{Selected Publications}

"Structure of the 30S Ribosomal Subunit" (Nature, 2000)

\clearpage

\section{Thomas A. Steitz}

\subsection*{Early Life and Education}

Thomas A. Steitz was born in 1940 in Milwaukee, Wisconsin, USA.

\subsection*{Career and Major Research}

Steitz earned bachelor's and doctoral degrees from Harvard University, in 1962 and 1966. He joined Yale University in 1970 and became a Sterling Professor of Molecular Biophysics and Biochemistry and an investigator of the Howard Hughes Medical Institute. With Peter Moore, he determined the atomic structure of the large (50S) subunit of the ribosome in 2000.

\subsection*{Nobel Prize-winning Work}

The work for which Thomas A. Steitz was awarded the Nobel Prize in Chemistry in 2009 was described by the Nobel Committee as: "for studies of the structure and function of the ribosome".

Steitz's structure of the 50S subunit showed where the peptide bond is formed and revealed the site at which many antibiotics stop protein synthesis.

\subsection*{Significance and Impact}

Because the large subunit carries out the chemistry of protein synthesis, its structure explained how proteins are made and how the ribosome can be blocked by drugs.

\subsection*{Later Career and Honors}

Steitz remained at Yale University for the rest of his career. He died on 2018-10-09 in Branford, Connecticut, USA.

\subsection*{Legacy}

His ribosome structures created molecular targets for a new generation of antibiotics.

\subsection*{Modern Developments}

Steitz's crystallographic structures of the ribosome revealed how the ribosome catalyzes protein synthesis and how antibiotics bind and block it. His work provided the atomic basis for understanding translation and for the rational design of antibiotics that treat infections such as tuberculosis, and it is foundational to structural biology.

\subsection*{Selected Publications}

"The Complete Atomic Structure of the Large Ribosomal Subunit at 2.4 Å Resolution" (Science, 2000)

\clearpage

\section{Ada E. Yonath}

\subsection*{Early Life and Education}

Ada E. Yonath was born in 1939 in Jerusalem, Israel.

\subsection*{Career and Major Research}

Yonath earned bachelor's, master's, and doctoral degrees at the Weizmann Institute of Science, completing her doctorate in 1968, followed by postdoctoral work at the Massachusetts Institute of Technology and the University of Chicago. She became a professor at the Weizmann Institute and director of its Kimmelman Center for Biomolecular Structure and Assembly.

\subsection*{Nobel Prize-winning Work}

The work for which Ada E. Yonath was awarded the Nobel Prize in Chemistry in 2009 was described by the Nobel Committee as: "for studies of the structure and function of the ribosome".

Yonath produced the first crystals of ribosomes from extremophilic bacteria that diffracted X-rays well, opening the way to the atomic structures determined by the prize.

\subsection*{Significance and Impact}

Her crystallization of the ribosome, a complex assembly of RNA and proteins, was a landmark that made high-resolution structure determination of the protein factory possible.

\subsection*{Later Career and Honors}

Yonath continues her research on the ribosome and its interaction with antibiotics at the Weizmann Institute of Science.

\subsection*{Legacy}

Her pioneering crystallization methods remain the basis for structural studies of the ribosome and of other large biological assemblies.

\subsection*{Modern Developments}

Yonath's determination of the three-dimensional structures of ribosomes, including those of antibiotic-resistant bacteria, revealed the molecular machinery of protein synthesis. Her work explains how antibiotics interfere with the ribosome and has guided the development of new antibacterial drugs to combat resistance.

\subsection*{Selected Publications}

"The Complete Atomic Structure of the Large Ribosomal Subunit at 2.4 Å Resolution" (Science, 2000)

\clearpage

\section*{Chapter Summary}

The 2009 prize recognized the determination of the structure and function of the ribosome, the cellular machine that translates the genetic code into proteins. The atomic structures of its subunits explain how proteins are made and how antibiotics act, and they have guided drug development.
